\section{Algorithms}\label{sec:algorithms}

The certificates of the preceding sections concern an exact operator. Using
them requires an algorithm that computes coordinate supports, decides when to
stop, and turns floating-point output into valid bounds.
This section describes two implementations with different contracts.

The first is an exact reference procedure for a declared rational relaxation
family. It commits an OBBT round only after every endpoint has been proved
optimal, and it stops with a fixed box only when the finite protected-box
certificate of Section~\ref{sec:certificates} holds. The second is a numerical
policy that runs inside a branch-and-bound solver. Each bound it applies is
justified by a dual bound evaluated with directed rounding, but its cutoff
comes from a numerically accepted incumbent, and its decisions about where to
spend auxiliary LPs are heuristic. Table~\ref{tab:contracts} summarizes the
difference; the
experiments of Section~\ref{sec:experiments} evaluate only the second
implementation.

\begin{table}[htbp]
\centering
\small
\begin{tabular}{@{}p{0.16\linewidth}p{0.38\linewidth}p{0.38\linewidth}@{}}
\toprule
& \raggedright Exact reference procedure
& \raggedright Numerical solver policy\tabularnewline
\midrule
\raggedright Relaxation
& \raggedright Fixed rational rows and McCormick product rows, rebuilt on each box
& \raggedright Outward-rounded local LP with a retained tangent pool, built at each callback\tabularnewline
\raggedright Endpoint proof
& \raggedright Exactly verified primal--dual pair
& \raggedright Dual bound with a residual term, evaluated with directed rounding\tabularnewline
\raggedright Cutoff
& \raggedright A given rational number $U$
& \raggedright A numerically accepted incumbent value, rounded upward\tabularnewline
\raggedright Stopping
& \raggedright Protected fixed box, round budget, or failed proof
& \raggedright Work budgets, triggers, screening, and a two-LP pilot\tabularnewline
\raggedright On failure
& \raggedright Inconclusive; earlier rounds are kept
& \raggedright No bound is proposed; native search continues\tabularnewline
\raggedright Guarantee
& \raggedright Every committed box and the stopping certificate
& \raggedright Each applied bound, conditional on the cutoff\tabularnewline
\bottomrule
\end{tabular}
\caption{The two implementations. Only the numerical policy was evaluated
in the computational comparison.}
\label{tab:contracts}
\end{table}

\subsection{The reference relaxation family}\label{sec:algorithms-family}

The exact procedure works with the reference family of
Section~\ref{sec:certificates} and rational data. A lifted point is
$z=(x,y)$, where $y_k$ stands for a product $x_{i_k}x_{j_k}$ and $i_k=j_k$ is
allowed. On a rational box $B=[\ell,u]$ the relaxation $R(B)$ consists of
fixed rational rows $Gz\le g$, the bound rows $\ell\le x\le u$, and the four
McCormick rows of every product; for a square these are the two endpoint
tangents and the secant. The objective is a rational affine function
$v(z)=c^\top z+c_0$, and $R_U(B)=\{z\in R(B):v(z)\le U\}$. By
Lemma~\ref{lem:reference-family}, the family is lifted monotone with a common
objective, and each $R(B)$ is a polytope. The infimum in the projected
objective~\eqref{eq:projected} is therefore attained whenever the fiber over
$x$ is nonempty, so that
\begin{equation}\label{eq:ref-projection}
 K_U(B)=\pi\bigl(R_U(B)\bigr),\qquad
 T_U(B)=\hull\pi\bigl(R_U(B)\bigr),
\end{equation}
where $\pi(z)=x$. The endpoints of $T_U(B)$ are thus the optimal values of the
$2n$ LPs that minimize and maximize each $x_i$ over $R_U(B)$.

Validity for an original problem is a premise that the caller must
establish. The McCormick rows supply valid lifts of the products, but the
fixed rows and the objective must be valid on the original sublevel set and
on the domain concerned. Equality of stored row data does not prove where a
row is valid.

\subsection{Exact endpoint certificates}\label{sec:algorithms-lp}

Each directional LP of the reference family has the form
$\min\{d^\top z:Az\le b\}$, where the rows $Az\le b$ list the fixed rows, the
bound rows, the McCormick rows, and the cutoff row, and $d=\pm e_i$ is a signed
unit vector. The following certificate is LP weak duality.

\begin{lemma}[Exact LP optimality certificate]\label{lem:lp-certificate}
Let $A\in\Q^{p\times N}$, $b\in\Q^p$, and $d\in\Q^N$. Suppose
$\bar z\in\Q^N$ and $\lambda\in\Q^p$ satisfy
\begin{equation}\label{eq:lp-certificate}
 A\bar z\le b,\qquad \lambda\le0,\qquad A^\top\lambda=d,\qquad
 d^\top\bar z=b^\top\lambda .
\end{equation}
Then $\bar z$ minimizes $d^\top z$ subject to $Az\le b$, and the optimal
value is $d^\top\bar z$.
\end{lemma}

\begin{proof}
Let $z$ satisfy $Az\le b$. Multiplying the $r$th inequality by
$\lambda_r\le0$ reverses it, so $\lambda^\top Az\ge\lambda^\top b$. Hence
$d^\top z=\lambda^\top Az\ge b^\top\lambda=d^\top\bar z$. Since $\bar z$ is
itself feasible, it is optimal.
\end{proof}

Every condition in~\eqref{eq:lp-certificate} is a finite set of rational
inequalities and identities. Checking them uses no tolerance and no solver
status. Because the bound rows are explicit rows of $A$, the multiplier
$\lambda$ also carries their dual information; the certificate therefore does
not depend on how a numerical solver treats variable bounds. A floating-point
LP solver supplies only a proposal $(\tilde z,\tilde\lambda)$. The reference
implementation replaces each proposed floating-point number by the rational
number with denominator at most $D=10^9$ that is closest to the number's
shortest decimal representation, and then checks~\eqref{eq:lp-certificate}
exactly. Checking floating-point proposals exactly after rational
reconstruction is an established technique of exact linear programming~\cite{applegate2007-exact-solutions-to-linear-programming}. A
failed check rejects the proposal. It says nothing about the LP itself.

The lemma needs an optimal solution. If $R_U(B)$ is empty, no certificate of
this form exists. The procedure below then stops without a conclusion. It
does not construct a Farkas certificate, so it never declares a box infeasible.

\subsection{Certified closure}\label{sec:algorithms-closure}

The certified closure procedure takes the family data, a rational box $B_0$,
a rational cutoff $U$, a round budget $K\ge0$, and a flag for an optional
objective proposal. It keeps a list of committed rounds and a counter of
proposal requests.

\begin{enumerate}
\item Set $B\leftarrow B_0$ and $k\leftarrow0$.
\item While $k<K$:
  \begin{enumerate}
  \item Build the rows of $R_U(B)$.
  \item For each $i=1,\ldots,n$, request proposals for
  $\min\{x_i:z\in R_U(B)\}$ and for $\min\{-x_i:z\in R_U(B)\}$, reconstruct
  rational values, and check~\eqref{eq:lp-certificate}. Call the verified
  optima $z^{(i,-)}$ and $z^{(i,+)}$. If any request returns no proposal, or
  any check fails, return \emph{inconclusive} with the current box $B$; the
  partial round is discarded.
  \item Commit the round: $B'=[\ell',u']$ with
  $\ell'_i=z^{(i,-)}_i$ and $u'_i=z^{(i,+)}_i$. Set $B\leftarrow B'$ and
  $k\leftarrow k+1$.
  \item Check exactly whether all $2n$ lifted points $z^{(i,\pm)}$, with every
  auxiliary coordinate, satisfy every row of $R_U(B')$, the relaxation
  rebuilt on the new box. If they do, optionally request a proposal for
  $\min\{v(z):z\in R_U(B')\}$; if its reconstruction satisfies every row of
  $R_U(B')$ exactly, add it to the points. Return \emph{fixed} with box $B'$
  and the lifted pool $\widehat W$ of checked points. Otherwise continue with
  the next round.
  \end{enumerate}
\item Return \emph{unfinished} with the current box $B$.
\end{enumerate}

Step~2(d) is the test of Proposition~\ref{prop:round-check}(a), applied to
the complete lifted optima of the round, which lie in $R_U(B)$ and whose
projections attain all endpoints of $B'$. The reported cost includes every
proposal request, including failed ones, and the elapsed time of all exact
checks. The optional objective proposal is never needed for a stopping
decision.

\begin{theorem}[Outcomes of certified closure]\label{thm:closure}
Run the procedure on the reference family with rational data. Write
$B_k$ for the box after $k$ committed rounds.
\begin{enumerate}
\item[(a)] \emph{Exact rounds.} Every committed round is the exact Jacobi
round: $B_{k+1}=T_U(B_k)\subseteq B_k$. Hence every returned box equals
$T_U^{\,k}(B_0)$, where $k$ is the number of committed rounds.
\item[(b)] \emph{Fixed status.} If the procedure returns \emph{fixed} with box
$P$ and lifted pool $\widehat W$, then $\pi(\widehat W)$ is a witness pool for
$P$ with threshold at most $U_{\widehat W}=\max_{z\in\widehat W}v(z)\le U$, and
$P$ is the greatest fixed box of $T_U$ in $B_0$ and the limit of the exact
Jacobi iterates. Every sequence of relaxation-sound steps that starts from a
box $C\supseteq P$ and uses cutoffs at least $U_{\widehat W}$ retains $P$, and
the relaxation bound of every box $C\supseteq P$ is at most
$\min_{z\in\widehat W}v(z)$.
\item[(c)] \emph{Prompt detection.} A committed round whose input box is a
fixed box of $T_U$ returns \emph{fixed}. Consequently every committed round
that does not return \emph{fixed} strictly shrinks the box, and if
$T_U(B_m)=B_m$, the budget allows $m+1$ rounds, and all proposals in these
rounds pass their checks, then the procedure returns \emph{fixed} after at
most $m+1$ rounds.
\item[(d)] \emph{Other statuses.} If the procedure returns \emph{unfinished}
or \emph{inconclusive}, the returned box is the last committed iterate. No
claim is made about the uncommitted round. In particular, a missing proposal
does not certify infeasibility.
\item[(e)] \emph{Validity.} If the family is valid for an original problem on
$B_0$, every returned box contains every original feasible point in $B_0$
with objective at most $U$.
\end{enumerate}
\end{theorem}

\begin{proof}
(a) By Lemma~\ref{lem:lp-certificate}, $z^{(i,-)}$ minimizes and $z^{(i,+)}$
maximizes $x_i$ over $R_U(B_k)$. By~\eqref{eq:ref-projection}, these optimal
values are the endpoints of $T_U(B_k)$, so the committed box is $T_U(B_k)$.
The bound rows of $R(B_k)$ give $B_{k+1}\subseteq B_k$. Induction over
committed rounds gives the second statement; failed rounds change nothing.

(b) Let $P=B_{k+1}=T_U(B_k)$ be the box just committed. The verified optima
$z^{(i,\pm)}$ lie in $R_U(B_k)$, and by part~(a) their projections attain all
endpoints of $P$. The check in step~2(d) is the hypothesis of
Proposition~\ref{prop:round-check}(a), which therefore makes these
projections a witness pool for $P$. An accepted objective point satisfies
every row of $R_U(P)$, including the bound rows, so its projection lies in
$K_U(P)$, and adding it leaves the hull equal to $P$. Every point $z$ of
$\widehat W$ lies in $R(P)$, so $\phi_P(\pi(z))\le v(z)$, which bounds the
threshold by $U_{\widehat W}$. Theorem~\ref{thm:protected} gives $T_U(P)=P$
and the persistence statement; since the reference family is defined on
every box, the theorem applies with any box containing $C$ in the role of
$B_0$. By Lemma~\ref{lem:greatest-fixed}, every Jacobi iterate contains the
greatest fixed box $P_U$ in $B_0$, so $P_U\subseteq P$, and $P\subseteq P_U$
because $P$ is a fixed box in $B_0$. All later Jacobi iterates equal $P$.
The bound on the relaxation value is Corollary~\ref{cor:objective-ceiling},
applied to the points of $\widehat W\subseteq R(P)$.

(c) If $T_U(B_k)=B_k$, the round commits $B_{k+1}=B_k$. The rows of
$R_U(B_{k+1})$ are then exactly those of $R_U(B_k)$, because the construction
depends only on the box and the fixed data, and the verified optima belong to
$R_U(B_k)$. The check therefore passes; this is
Proposition~\ref{prop:round-check}(b). If a committed round does not return
\emph{fixed}, its input box is therefore not fixed, and
$B_{k+1}=T_U(B_k)\ne B_k$. For the last statement, either the check succeeds
earlier, or round $m+1$ starts at the fixed box $B_m$.

(d) This follows from (a): a round changes the box only when it is committed.

(e) By Lemma~\ref{lem:order}(a), each original feasible point $x\in B_k$ with
objective at most $U$ lies in $T_U(B_k)=B_{k+1}$. Induction from $B_0$ proves
the claim.
\end{proof}

Part~(b) states what the stopping certificate guarantees: the returned box
is the greatest fixed box and the limit of exact Jacobi iteration at cutoff
$U$, and no order of exact or conservative coordinate updates in this family,
at cutoffs at least $U_{\widehat W}$, can remove any part of it. Because
$U_{\widehat W}$ can lie below $U$, the guarantee survives incumbent
improvements down to that value. The relaxation-value ceiling bounds the
possible improvement of this
relaxation's lower bound; it does not bound the gain from cuts, integrality
reasoning, or a different relaxation. Part~(c) shows that the check after a
round is at most one round late. It can be late: the verified optima of a
round can violate a row rebuilt on the new box although that box is fixed, as
in Example~\ref{ex:round-check}, and only the next round, which starts at the
fixed box, produces points that pass.

Every run stops, because the budget $K$ is finite, but no finite budget
guarantees the status \emph{fixed}. That status requires exact iterates that
reach a fixed box, and the square examples below never do. Rational
reconstruction from double-precision proposals can also fail on
well-behaved LPs once the exact iterates need many digits. The procedure
then reports \emph{inconclusive} and keeps the last committed box. An exact
rational LP solver could replace
the floating-point proposals without changing the checks. It would remove
reconstruction failures, but not the growth of the iterates: in Heron's
iteration below, the number of digits roughly doubles in every round.

Table~\ref{tab:closure-runs} records the archived outcomes of the reference
implementation on small fixtures. They illustrate the three statuses; they are
correctness checks, not performance evidence. The first fixture is
Example~\ref{ex:three-variable}: the cube is fixed at cutoff $0$, so the first
round starts at a fixed box, and the objective proposal certifies the ceiling
$-3$, the relaxation bound, while the original optimum is $0$. The linear
fixture has the rows $x_1+x_2\ge1$ and $x_1=x_2$ on $[0,2]^2$. In the square
fixtures, the cutoff $0$ halves the upper endpoint in every round, as in
Example~\ref{ex:halving}, and never reaches the fixed box $\{0\}$. With cutoff
$1$ the upper endpoint follows Heron's iteration of Example~\ref{ex:heron},
$2$, $5/4$, $41/40$, $3281/3280$; the next iterate is $21523361/21523360$, and
the denominators roughly square in every round. In the fourth round, the
reconstructed proposal for the upper direction violated a row exactly, and the
procedure returned \emph{inconclusive} with the third box.

\begin{table}[htbp]
\centering
\small
\begin{tabular}{@{}lrrll@{}}
\toprule
Fixture & Committed rounds & Proposals & Status & Returned box\\
\midrule
Three-variable stall, cutoff $0$ & 1 & 7 & fixed & $[-1,1]^3$, ceiling $-3$\\
Linear coupling & 1 & 4 & fixed & $[1/2,2]^2$\\
Interval with negative bounds & 1 & 2 & fixed & $[-3,-1]$\\
Square, cutoff $0$, budget 5 & 5 & 10 & unfinished & $[0,1/32]$\\
Square, cutoff $1$, budget 3 & 3 & 6 & unfinished & $[0,3281/3280]$\\
Square, cutoff $1$, budget 12 & 3 & 8 & inconclusive & $[0,3281/3280]$\\
\bottomrule
\end{tabular}
\caption{Archived runs of the certified closure procedure. The square fixtures
start at $[0,1]$ (cutoff $0$) and $[0,2]$ (cutoff $1$). The stall count includes
one objective proposal. Further archived cases inject an invalid dual vector,
a missing proposal after one committed round, and an invalid objective
proposal; they return \emph{inconclusive} with no committed round,
\emph{inconclusive} with the first committed box, and \emph{fixed} with the
certificate unchanged, respectively. A zero round budget returns
\emph{unfinished} without a proposal.}
\label{tab:closure-runs}
\end{table}

\subsection{A numerical policy inside branch and bound}\label{sec:algorithms-policy}

The second implementation adds a propagator to SCIP for quadratically
constrained problems with quadratic objectives. It leaves SCIP's own
propagation, including its LP-based
OBBT~\cite{gleixner2017-three-enhancements-for-optimization-based},
unchanged. At selected nodes it
builds a separate LP relaxation of the current node box, solves some
directional LPs with a floating-point solver, and proposes valid local bounds.
The time that the propagator records for this work is its \emph{propagator
time}. The propagator never uses an LP status to prune a node.

\subsubsection{Local relaxation}

Let $B=[\ell,u]$ be the current local box of the original variables, read
from SCIP and intersected with the original bounds; every bound must be
finite, otherwise the callback records an unsupported box and returns. The
product terms are the distinct pairs $(i,j)$, $i\le j$, that occur in the
objective or a row. The lifted LP has one variable $y_{ij}$ per term and these
rows:
\begin{itemize}
\item the four McCormick rows~\eqref{eq:mccormick-gap} for each $i\ne j$;
\item for each square, tangent rows $y_{ii}\ge2px_i-p^2$ at five equally spaced
points $p$ of $[\ell_i,u_i]$ and at retained points from earlier callbacks, and
the secant $y_{ii}\le(\ell_i+u_i)x_i-\ell_iu_i$;
\item every original row, with each product replaced by its lifted variable;
\item if a cutoff $U_c$ is available, the objective in minimization form, with
products replaced, bounded by $U_c$.
\end{itemize}
The variable bounds are $x\in B$ and, for each term, an interval containing
every product value on $B$: the four endpoint products for $i\ne j$, and
$[\min_{s\in[\ell_i,u_i]}s^2,\max\{\ell_i^2,u_i^2\}]$ for a square. These
intervals are computed in floating point and widened by one unit in the last
place in each direction. A tangent is valid for the square everywhere, so a
retained point outside the current interval still yields a valid row. At
every callback the five points of the current interval are offered to a
pool shared by all nodes, which keeps the first $25$ distinct points of
each square and accepts no more.

The relaxation is frozen during a callback. Bounds applied in that callback
change SCIP's domains but not the auxiliary LP, so a callback performs part
of a frozen round; the next callback rebuilds. The relaxations of different
callbacks need not form a monotone family, for two reasons. First, the five
tangent points move with the box, and such box-relative points violate
monotonicity on their own: for $f(x)=x^2$ these rows alone give
$T_{1/4}([-1,1])=[-1/2,1/2]$ but $T_{1/4}([-1/2,1])=[-1/2,41/80]$
(Section~\ref{sec:foundations}). The pool restores the tangents of earlier
boxes only while it has room. Second, tangents that the pool gains later are
rows that were absent when lifted points were checked on an earlier box, and
they can cut those points. The future-round results therefore do not apply
to this sequence of relaxations.

\subsubsection{Valid bounds from floating-point data}

The model data are stored in binary floating point. The implementation treats
each stored number as the exact dyadic rational it represents, constructs every
new row exactly, and then rounds its coefficients to floating point with a
correction of the right-hand side.

Each candidate uses the dual bound of Proposition~\ref{prop:dual-residual},
evaluated with directed rounding as justified by Lemma~\ref{lem:enclosure},
for an LP whose rows are rounded outward using
Proposition~\ref{prop:row-correction}. Appendix~\ref{sec:numerical-validation}
gives the formulas and proofs. The argument assumes IEEE binary64 arithmetic
with gradual underflow and with flush-to-zero and denormals-are-zero modes
disabled. It proves validity for the stored binary model on the local box,
without asserting equivalence to a decimal or symbolic model before import.

For a lower direction the solver minimizes $x_i$ and the candidate bound is
$\tilde\mu$; for an upper direction it minimizes $-x_i$ and the candidate is
$-\tilde\mu$. Each candidate is moved outward by a further margin
$10^{-7}\max\{1,|\tilde\mu|\}$ and rounded outward. A candidate is proposed to
SCIP only if it lies in the current local domain, improves the current bound by
more than $\max\{10^{-9},10^{-4}(u_i-\ell_i)\}$, and, for an integer variable,
remains inside the domain after rounding to the next integer. A candidate that
would empty the domain is discarded rather than used to prune the node.
Accepted bounds are local to the node; SCIP restores them on backtracking.
LP errors, time limits, infeasibility reports, and unbounded statuses produce
no candidate.

\begin{proposition}[Conditional validity of applied bounds]\label{prop:sidecar-validity}
Let $x^\star\in B$ satisfy the stored original rows, bounds, and integrality
restrictions. If the callback uses a cutoff $U_c$, assume also that the stored
objective, in minimization form, satisfies $f(x^\star)\le U_c$. Then $x^\star$
satisfies every bound that the callback applies.
\end{proposition}

\begin{proof}
The exact lift of $x^\star$ satisfies the McCormick, tangent, and secant rows,
which are valid inequalities for the products on $B$, the product intervals,
the lifted original rows, and the cutoff row. By
Proposition~\ref{prop:row-correction} it satisfies every rounded row, so it is
feasible for the stored LP. Hence $x^\star_i\ge\mu\ge\tilde\mu$ for a lower
direction and $x^\star_i\le-\tilde\mu$ for an upper direction. The outward
margin weakens each candidate further. SCIP rounds an integer bound to the
next integer, which an integer $x^\star_i$ still satisfies.
\end{proof}

The cutoff is the weak point of this guarantee. It is computed from SCIP's best
solution $\hat x$ only if $\hat x$ passes a numerical check of all original
rows, bounds, and integrality with absolute tolerance $10^{-6}$. Its objective
is then evaluated exactly in rational arithmetic, rounded upward, and increased
by $10^{-6}\max\{1,|f(\hat x)|\}$. Exact evaluation prevents cancellation; for
example, $10^{16}x^2+xy-10^{16}y^2$ at $x=y=1$ has value $1$, while sequential
floating-point summation can return $0$. Thus $U_c\ge f(\hat x)$. If $\hat x$
is exactly feasible, then every global minimizer of the stored model satisfies
$f\le U_c$, and Proposition~\ref{prop:sidecar-validity} shows that no applied
bound removes one. If $\hat x$ is only approximately feasible, its value may lie
below the optimal value, and the conclusion can fail. Without an incumbent no
cutoff row is used, and no feasible point of the node is removed. This is the
ordinary numerical contract of the solver, made explicit; neither the
propagator nor the solve as a whole is an exact certificate.

\subsubsection{Scheduling policy}

The fixed and adaptive variants share the relaxation, the bound validation,
the budgets, and the triggers. They differ in direction order, witness
screening, and a pilot rule. Table~\ref{tab:policy-parameters} lists the
parameters, which were frozen before the comparison of
Section~\ref{sec:experiments}. SCIP calls the propagator at every node before
the node LP, and it may call it several times at a node.

\begin{table}[htbp]
\centering
\small
\begin{tabular}{@{}lr@{}}
\toprule
Parameter & Value\\
\midrule
Directional LPs per run & 64\\
Directional LPs per callback at the root / elsewhere & 12 / 4\\
Acting callbacks per run; per node & 24; 3\\
Maximum depth & 64\\
Propagator time allowance & 10\% of the time limit\\
Time limit of one directional LP & 0.25 s\\
Incumbent retrigger: decrease of the cutoff by at least & $10^{-2}\max\{1,|U_{\mathrm{old}}|\}$\\
Domain retrigger: relative width decrease of some variable & $0.10$\\
Minimum improvement of a proposed bound & $\max\{10^{-9},10^{-4}(u_i-\ell_i)\}$\\
Outward margin of a proposed bound & $10^{-7}\max\{1,|\tilde\mu|\}$\\
Screening tolerance (adaptive) & $10^{-7}\max\{1,u_i-\ell_i\}$\\
Pilot (adaptive): LPs; minimum total normalized gain & 2; $10^{-3}$\\
Cached LP points (adaptive); retained tangent points per square & 24; 25\\
\bottomrule
\end{tabular}
\caption{Frozen parameters of the numerical policy.}
\label{tab:policy-parameters}
\end{table}

\emph{Admission.} A call does nothing if the run has used all its directional
LPs or reached its limit of acting callbacks, defined below, if the
propagator time allowance or the solve deadline is exhausted, or if the
node depth exceeds the limit. These checks are cheap and are not timed;
Section~\ref{sec:effort} accounts for them separately.

\emph{Triggers.} After admission the propagator's timer starts. The
callback reads the local box and computes the cutoff, and it continues only
if a trigger
fires. A propagation call that passes admission and a trigger is an
\emph{acting callback}. Only acting callbacks build a relaxation, and the
callback limits
of Table~\ref{tab:policy-parameters} count them. For each node number the
policy records the box and the cutoff of the last acting visit and the
number of acting callbacks at the node. A first visit triggers. A later
visit triggers only if fewer than three acting callbacks have occurred at
the node and either the cutoff decreased by
at least $10^{-2}\max\{1,|U_{\mathrm{old}}|\}$ (or became finite), or some
variable's width decreased by at least $10\%$ relative to the recorded box.
A visit that triggers nothing returns, but its time is charged.
For a monotone relaxation family, a smaller box or a smaller cutoff shrinks
the cutoff sets (Lemma~\ref{lem:order}) and can move supports that were
optimal before, and the cutoff-dependent floors of
Section~\ref{sec:local-rates} motivate the incumbent trigger. The policy's
own relaxation is not monotone, as explained above, so this reasoning
motivates the two triggers without justifying them. The policy does not
revisit directions left unprocessed at an unchanged node, and the
thresholds are heuristic.

\emph{Directions and limits.} Candidate variables are those occurring in a
product term whose width exceeds $10^{-9}$; each contributes its lower
direction followed by its upper direction. The fixed variant uses variable
index order. The adaptive variant orders variables by the score
\[
 \mathrm{score}_i=\sum_{(i,j)}\omega_{ij}(u_i-\ell_i)(u_j-\ell_j),
\]
summed over the product terms containing $i$, where $\omega_{ij}$ is the sum
of the absolute coefficients of the term in the objective and all rows. A
product of factor widths measures the scale of the term's McCormick error on
the box, so the score favors variables whose relaxation error is large. Each
callback may solve up to 12 directional LPs at the root and 4 elsewhere,
limited by the LPs remaining in the run. Each LP receives the time limit
$\min\{0.25\,\mathrm{s},\ \text{remaining allowance and solve time}\}$, raised
to at least one millisecond.

\emph{Screening (adaptive).} The adaptive variant keeps the 24 most recent
LP solutions from any node. At a callback it first checks each cached point
for exact feasibility in the current frozen LP, in rational arithmetic after a
fast interval test. A point that passes and lies within
$10^{-7}\max\{1,u_i-\ell_i\}$ of an endpoint screens that direction. This is the
current-round ceiling of Proposition~\ref{prop:current-round}, with a
tolerance: the screened direction can gain at most that tolerance in this
frozen LP. After every solved LP the new solution is checked and used in the
same way. The check is repeated at every callback, because a point from another
node, a larger box, or a weaker cutoff need not satisfy the current rows.

\emph{Pilot (adaptive).} After the second LP of a callback, the adaptive
variant stops if the gains of the bounds that SCIP accepted in this callback
sum to less than $10^{-3}$. The gain of a bound is the improvement of the
proposed value over the current bound, measured before SCIP rounds integer
bounds, divided by the variable's width at the start of the callback. The
accumulated sum never decreases, so this is a single test. Without screening,
the first two LPs are the two directions of the highest-ranked variable, and
the pilot asks whether that variable can be tightened noticeably. It is a
heuristic: it does not bound the gain in unsampled directions or in later
rounds.

\emph{Timing.} A callback stops when its LP limit or its share of the
remaining propagator allowance is used up. Construction checks the deadline every
64 rows, and screening checks it between points. Individual construction steps,
LP calls, and validations are not interrupted, and the one-millisecond floor
of the LP limit can exceed the remaining allowance, so a callback can overrun
its deadline; the overrun is charged.

\subsubsection{What the policy uses from the theory}

The policy uses three proved statements: the outward correction and the
residual dual bound, which make each applied bound valid under the cutoff
premise; and the current-round ceiling of
Proposition~\ref{prop:current-round}, which makes screening sound for the
frozen LP. It does not use protected boxes, cutoff thresholds of witness
pools, or residual tail bounds to stop tightening permanently, and it neither
searches for nor checks such certificates. Using them would require changes
that the policy does not implement. Its relaxations are not monotone, so a
box protected for the relaxation of one callback need not remain protected
for the next. A protected box also covers only steps at cutoffs at least its
threshold, whereas incumbent improvements lower the cutoff during the
search. SCIP's own propagation, branching,
and integer rounding shrink the node box by operations other than the
supports of the auxiliary LP, so a box protected against the propagator's
operator can still be cut,
and its containment would have to be rechecked at every node. Finally,
proposing a protected box from support LPs takes a complete round. Every
model of the prospective comparison has at least ten variables in product
terms, so a complete round over them takes at least 20 LPs, more than the 12
allowed per root callback, unless earlier bounds have fixed some of these
variables. The measured comparison is
therefore a test of this policy and its costs, not of the all-future
certificates.
